\documentclass[10pt,a4paper]{amsart}
\usepackage{amsmath,amssymb,amsthm,graphicx}
% Retained minimal preamble for the Stacks Project chapter sites-cohomology.tex (Cohomology on Sites).
% Verbatim excerpt of the chapter preamble: Xy-pic setup (native preamble lines 14-18) and the native
% theorem declarations with the mathematical macros (native preamble lines 154-217). Copied mechanically;
% nothing added, removed or redefined. The wrapper class and the amsmath/amssymb/amsthm/graphicx base
% packages are supplied by the assembler, so the class/externaldocument part of the native preamble is omitted.
% For commutative diagrams we use Xy-pic
\usepackage[all]{xy}

% We use 2cell for 2-commutative diagrams.
\xyoption{2cell}
% Theorem environments.
%
\theoremstyle{plain}
\newtheorem{theorem}[subsection]{Theorem}
\newtheorem{proposition}[subsection]{Proposition}
\newtheorem{lemma}[subsection]{Lemma}

\theoremstyle{definition}
\newtheorem{definition}[subsection]{Definition}
\newtheorem{example}[subsection]{Example}
\newtheorem{exercise}[subsection]{Exercise}
\newtheorem{situation}[subsection]{Situation}

\theoremstyle{remark}
\newtheorem{remark}[subsection]{Remark}
\newtheorem{remarks}[subsection]{Remarks}

\numberwithin{equation}{subsection}

% Macros
%
\def\lim{\mathop{\mathrm{lim}}\nolimits}
\def\colim{\mathop{\mathrm{colim}}\nolimits}
\def\Spec{\mathop{\mathrm{Spec}}}
\def\Hom{\mathop{\mathrm{Hom}}\nolimits}
\def\Ext{\mathop{\mathrm{Ext}}\nolimits}
\def\SheafHom{\mathop{\mathcal{H}\!\mathit{om}}\nolimits}
\def\SheafExt{\mathop{\mathcal{E}\!\mathit{xt}}\nolimits}
\def\Sch{\mathit{Sch}}
\def\Mor{\mathop{\mathrm{Mor}}\nolimits}
\def\Ob{\mathop{\mathrm{Ob}}\nolimits}
\def\Sh{\mathop{\mathit{Sh}}\nolimits}
\def\NL{\mathop{N\!L}\nolimits}
\def\CH{\mathop{\mathrm{CH}}\nolimits}
\def\proetale{{pro\text{-}\acute{e}tale}}
\def\etale{{\acute{e}tale}}
\def\QCoh{\mathit{QCoh}}
\def\Ker{\mathop{\mathrm{Ker}}}
\def\Im{\mathop{\mathrm{Im}}}
\def\Coker{\mathop{\mathrm{Coker}}}
\def\Coim{\mathop{\mathrm{Coim}}}

% Boxtimes
%
\DeclareMathSymbol{\boxtimes}{\mathbin}{AMSa}{"02}

%
% Macros for moduli stacks/spaces
%
\def\QCohstack{\mathcal{QC}\!\mathit{oh}}
\def\Cohstack{\mathcal{C}\!\mathit{oh}}
\def\Spacesstack{\mathcal{S}\!\mathit{paces}}
\def\Quotfunctor{\mathrm{Quot}}
\def\Hilbfunctor{\mathrm{Hilb}}
\def\Curvesstack{\mathcal{C}\!\mathit{urves}}
\def\Polarizedstack{\mathcal{P}\!\mathit{olarized}}
\def\Complexesstack{\mathcal{C}\!\mathit{omplexes}}
% \Pic is the operator that assigns to X its picard group, usage \Pic(X)
% \Picardstack_{X/B} denotes the Picard stack of X over B
% \Picardfunctor_{X/B} denotes the Picard functor of X over B
\def\Pic{\mathop{\mathrm{Pic}}\nolimits}
\def\Picardstack{\mathcal{P}\!\mathit{ic}}
\def\Picardfunctor{\mathrm{Pic}}
\def\Deformationcategory{\mathcal{D}\!\mathit{ef}}
\usepackage{cleveref}
\crefname{originalproposition}{original proposition}{original propositions}
\Crefname{originalproposition}{Original proposition}{Original propositions}
\makeatletter
\newlabel{derived-proposition-left-derived-exists}{{0794}{source}{}{sourcepointer-1}{}}
\newlabel{derived-proposition-left-derived-exists@cref}{{[originalproposition][0][]0794}{[1][0][]source}{}{}{}}
\makeatother
\crefname{originallemma}{original lemma}{original lemmas}
\Crefname{originallemma}{Original lemma}{Original lemmas}
\makeatletter
\newlabel{lemma-compute-by-cosimplicial-resolution}{{08Q9}{source}{}{sourcepointer-2}{}}
\newlabel{lemma-compute-by-cosimplicial-resolution@cref}{{[originallemma][0][]08Q9}{[1][0][]source}{}{}{}}
\makeatother
\crefname{originallemma}{original lemma}{original lemmas}
\Crefname{originallemma}{Original lemma}{Original lemmas}
\makeatletter
\newlabel{sites-modules-lemma-module-quotient-flat}{{03EW}{source}{}{sourcepointer-3}{}}
\newlabel{sites-modules-lemma-module-quotient-flat@cref}{{[originallemma][0][]03EW}{[1][0][]source}{}{}{}}
\makeatother
\title{Under a cosimplicial resolution and an isomorphism of L pi\_! on the structure sheaves, L pi\_! K agrees with L pi\_! of the extended scalars}
\author{The Stacks Project Authors}
\date{stacks-project:sites-cohomology.tex:local-snapshot-20261011; GNU FDL 1.2 or later}
\begin{document}
\maketitle
\noindent\textit{Editorial note.} All delivered mathematical text is the author's own, copied from the pinned Stacks Project chapter file sites-cohomology.tex; no statement or proof step was rewritten, reordered or removed.\par\medskip
\noindent\textit{Editorial note.} The counted proof is the author's complete argument: the proof is written out in full: the reduction to bounded above complexes and to the flat sheaves j\_\{U!\}O\_U, the definition of the simplicial abelian groups X\_bullet and X'\_bullet, the functoriality argument giving the homotopy equivalences with P\_bullet and P'\_bullet, and the conclusion from the assumed quasi-isomorphism.\par\medskip
\noindent\textit{Editorial note.} The label lemma-O-homology-qis names the counted statement, and the proof environment immediately following it in the source is delivered verbatim as its proof.\par\medskip
\noindent\textit{Editorial note.} The statement and proof are copied from the pinned chapter file; the three cited inputs are kept as external pointers with their canonical Stacks tags.\par\medskip
\noindent\textit{Editorial note.} The preamble is the chapter's own retained minimal preamble (sources/stacks-sites-cohomology/retained-preamble.tex), a verbatim excerpt of the native preamble: the Xy-pic setup and the native theorem declarations with the mathematical macros.\par\medskip
\noindent\textit{Allowed standard prerequisites (editorial).}\begin{itemize}
\item The functor L pi\_! (homology of a site) and its computation by a cosimplicial resolution. Same chapter, the left derived functor of pi\_! and the lemma computing it on j\_\{U!\}O\_U through the morphism complex of a cosimplicial object U\_bullet. Supplies the identification of the complex X\_bullet with L pi\_!(j\_\{U!\}O\_U) and of P\_bullet, P'\_bullet with L pi\_! O, L pi\_! O'.
\item The chaotic topology and the flatness of j\_\{U!\}O\_U. Same chapter, the structure of sheaves on a category with chaotic topology: every bounded above complex of O-modules is quasi-isomorphic to one whose terms are direct sums of j\_\{U!\}O\_U. Supplies the reduction of the computation to the single sheaves j\_\{U!\}O\_U, together with their flatness, which makes tensor with O' exact.
\end{itemize}\medskip
\ifdefined\hypertarget\hypertarget{sourcepointer-1}{}\fi
\noindent\textit{Source reference.} Original proposition 0794: sources/stacks-chow/tags.txt, source lines 9259--9259. Prerequisite pointer invoked at proof line(s) 10444; it names a standard fact of a sibling chapter that the counted argument uses, and every step of the counted proof built on it is retained verbatim in the delivered body.\par\medskip
\ifdefined\hypertarget\hypertarget{sourcepointer-2}{}\fi
\noindent\textit{Source reference.} Original lemma 08Q9: sources/stacks-sites-cohomology/pinned-original-source.tex, source lines 10229--10229. Prerequisite pointer invoked at proof line(s) 10428; it names a standard fact of the same chapter that the counted argument uses, and every step of the counted proof built on it is retained verbatim in the delivered body.\par\medskip
\ifdefined\hypertarget\hypertarget{sourcepointer-3}{}\fi
\noindent\textit{Source reference.} Original lemma 03EW: sources/stacks-chow/tags.txt, source lines 4471--4471. Prerequisite pointer invoked at proof line(s) 10449; it names a standard fact of a sibling chapter that the counted argument uses, and every step of the counted proof built on it is retained verbatim in the delivered body.\par\medskip
\begin{lemma}
\label{lemma-O-homology-qis}
Let $\mathcal{C}$ be a category (endowed with chaotic topology).
Let $\mathcal{O} \to \mathcal{O}'$ be a map of sheaves of rings on
$\mathcal{C}$. Assume
\begin{enumerate}
\item there exists a cosimplicial object $U_\bullet$ in $\mathcal{C}$
as in Lemma \ref{lemma-compute-by-cosimplicial-resolution}, and
\item $L\pi_!\mathcal{O} \to L\pi_!\mathcal{O}'$ is an isomorphism.
\end{enumerate}
For $K$ in $D(\mathcal{O})$ we have
$$
L\pi_!(K) = L\pi_!(K \otimes_\mathcal{O}^\mathbf{L} \mathcal{O}')
$$
in $D(\textit{Ab})$.
\end{lemma}

\begin{proof}
Note: in this proof $L\pi_!$ denotes the left derived functor
of $\pi_!$ on abelian sheaves.
Since $L\pi_!$ commutes with colimits, it suffices
to prove this for bounded above complexes of $\mathcal{O}$-modules
(compare with argument of
Derived Categories, Proposition \ref{derived-proposition-left-derived-exists}
or just stick to bounded above complexes).
Every such complex is quasi-isomorphic to a bounded above complex
whose terms are direct sums of $j_{U!}\mathcal{O}_U$ with
$U \in \Ob(\mathcal{C})$, see
Modules on Sites, Lemma \ref{sites-modules-lemma-module-quotient-flat}.
Thus it suffices to prove the lemma
for $j_{U!}\mathcal{O}_U$. By assumption
$$
S_\bullet = \Mor_\mathcal{C}(U_\bullet, U)
$$
is a simplicial set homotopy equivalent to the constant simplicial
set on a singleton. Set $P_n = \mathcal{O}(U_n)$ and
$P'_n = \mathcal{O}'(U_n)$. Observe that the complex associated to the
simplicial abelian group
$$
X_\bullet : n \longmapsto \bigoplus\nolimits_{s \in S_n} P_n
$$
computes $L\pi_!(j_{U!}\mathcal{O}_U)$ by
Lemma \ref{lemma-compute-by-cosimplicial-resolution}.
Since $j_{U!}\mathcal{O}_U$ is a flat $\mathcal{O}$-module we have
$j_{U!}\mathcal{O}_U \otimes^\mathbf{L}_\mathcal{O} \mathcal{O}' =
j_{U!}\mathcal{O}'_U$ and $L\pi_!$ of this is computed by the complex
associated to the simplicial abelian group
$$
X'_\bullet : n \longmapsto \bigoplus\nolimits_{s \in S_n} P'_n
$$
As the rule which to a simplicial set $T_\bullet$ associated the simplicial
abelian group with terms $\bigoplus_{t \in T_n} P_n$ is a functor, we see
that $X_\bullet \to P_\bullet$ is a homotopy equivalence of simplicial
abelian groups. Similarly, the rule which to a simplicial set
$T_\bullet$ associates the simplicial abelian group with terms
$\bigoplus_{t \in T_n} P'_n$ is a functor. Hence $X'_\bullet \to P'_\bullet$
is a homotopy equivalence of simplicial abelian groups.
By assumption $P_\bullet \to P'_\bullet$ is a quasi-isomorphism
(since $P_\bullet$, resp.\ $P'_\bullet$ computes $L\pi_!\mathcal{O}$,
resp. $L\pi_!\mathcal{O}'$ by
Lemma \ref{lemma-compute-by-cosimplicial-resolution}).
We conclude that $X_\bullet$ and $X'_\bullet$ are quasi-isomorphic as desired.
\end{proof}

\end{document}
